\documentclass[11pt,a4paper]{article}
\usepackage[utf8]{inputenc}
\usepackage{amsmath,amsthm,amssymb}
\usepackage{mathtools}
\usepackage{hyperref}
\usepackage{cleveref}
\usepackage{geometry}
\geometry{margin=1in}

% Theorem environments
\newtheorem{theorem}{Theorem}[section]
\newtheorem{lemma}[theorem]{Lemma}
\newtheorem{proposition}[theorem]{Proposition}
\newtheorem{corollary}[theorem]{Corollary}
\newtheorem{conjecture}[theorem]{Conjecture}
\theoremstyle{definition}
\newtheorem{definition}[theorem]{Definition}
\newtheorem{example}[theorem]{Example}
\theoremstyle{remark}
\newtheorem{remark}[theorem]{Remark}

\title{Locality Tradeoffs for 2D Quantum LDPC Codes: An Expository Restatement of the BPT Bound with Range Bookkeeping}
\author{David}
\date{\today}

\begin{document}

\maketitle

\begin{abstract}
We study tradeoffs between rate, distance, and geometric locality for 2D quantum LDPC stabilizer codes. The main technical input is the 2D Bravyi--Poulin--Terhal (BPT) bound $k d^2 = O(n)$ for geometrically local commuting constraints; our contribution is expository and focuses on stating the assumptions cleanly and tracking how the hidden constant can depend on an explicit interaction range via a coarse-graining reduction to constant-range interactions. We also contrast these locality limitations with recent asymptotically good (necessarily non-geometrically-local) quantum LDPC constructions, and we explain how separator phenomena obstruct expander-like behavior in bounded-dimensional locality graphs.
\end{abstract}

\section{Introduction}

Quantum error correction is fundamental to fault-tolerant quantum computation. Among the various code families, quantum low-density parity-check (LDPC) codes have emerged as particularly promising candidates due to their favorable scaling properties and potential for efficient decoding.

The central question in quantum LDPC code theory concerns the tradeoffs between code parameters:
\begin{itemize}
    \item Rate $R = k/n$ (logical qubits per physical qubit)
    \item Distance $d$ (number of errors the code can correct)
    \item Maximum check weight $w$ (sparsity of parity checks)
    \item Geometric locality constraints
\end{itemize}

Classical results such as the Bravyi--Poulin--Terhal (BPT) bound \cite{BPT} establish fundamental limits for geometrically local codes. This note focuses on the 2D setting and keeps track of how the hidden constant in the $O(n)$ notation can depend on an explicit interaction range parameter (via coarse-graining). We do \emph{not} claim a new asymptotic tradeoff beyond BPT; rather, we aim for a clean, referee-proof statement of assumptions and bookkeeping.

\subsection{Main Contributions}

In this note, we:
\begin{enumerate}
    \item Restate the 2D BPT locality tradeoff with an explicit interaction range parameter using a coarse-graining reduction
    \item Explain (at a high level) why expander-based LDPC mechanisms are incompatible with fixed-dimensional geometric locality, via separator phenomena
    \item Demonstrate 2D tightness up to constants via the toric code
\end{enumerate}

\section{Preliminaries}

\subsection{Quantum LDPC Codes}

\begin{definition}[CSS Code Construction]
A Calderbank-Shor-Steane (CSS) code is a particular type of quantum error-correcting code constructed from two classical linear block codes, $C_1$ and $C_2$, with parity check matrices $H_X$ and $H_Z$, such that $C_2^\perp \subseteq C_1$. This condition is equivalent to the orthogonality constraint $H_X H_Z^T = 0$. The resulting quantum code corrects $X$-errors using $H_Z$ and $Z$-errors using $H_X$.
\end{definition}

\begin{definition}[Quantum LDPC Code]
A quantum LDPC code is a CSS code $\mathcal{C} = [[n, k, d]]$ defined by sparse parity-check matrices $H_X$ and $H_Z$ satisfying $H_X H_Z^T = 0$, with bounded stabilizer weight and bounded qubit degree. Concretely, suppose each row of $H_X$ and $H_Z$ has Hamming weight at most $w=O(1)$ (maximum stabilizer/check weight) and each column has Hamming weight at most $q=O(1)$ (maximum number of checks acting on a qubit). Both bounds are independent of the block length $n$.
\end{definition}

\begin{definition}[Tanner Graph]
The Tanner graph $G = (V \cup C, E)$ of a quantum LDPC code has variable nodes $V$ corresponding to qubits and check nodes $C$ corresponding to stabilizer generators, with edges indicating non-trivial support.
\end{definition}

\subsection{Geometric Locality}

\begin{definition}[Geometric Locality (bounded density, finite range)]
A quantum stabilizer code on $n$ qubits is \emph{geometrically local} in $\mathbb{R}^D$ with interaction range $L$ and density $\rho$ if:
\begin{enumerate}
    \item There exists an embedding of the $n$ qubits into $\mathbb{R}^D$ such that every ball of unit volume contains at most $\rho$ qubits.
    \item Each stabilizer generator acts only on qubits contained within some ball of diameter at most $L$.
\end{enumerate}
For asymptotic analysis, we typically assume $\rho = O(1)$ and define locality by the constraint that $L$ is constant independent of $n$.
\end{definition}

\begin{remark}[Bookkeeping for $L$ and $w$]
In a bounded-density embedding, a ball of diameter $L$ contains at most $O(\rho L^D)$ qubits. Thus, in many geometric models one has an implicit bound $w \le O(\rho L^D)$ for strictly local stabilizer generators. In this note we keep $w$ and $L$ as explicit parameters to cleanly separate ``range'' from ``weight,'' but one should keep in mind that they may not be independent in a given hardware model.
\end{remark}

\subsection{Expander Graphs and Spectral Methods}

\begin{definition}[Spectral Expansion]
A graph $G$ is a $\lambda$-expander if the second largest eigenvalue of its normalized adjacency matrix satisfies $|\lambda_2| \leq \lambda$.
\end{definition}

\begin{lemma}[Expander Mixing Lemma]
For a $\lambda$-expander $G$ with $n$ vertices and (regular) degree $r$, and any two vertex sets $S, T \subseteq V$:
\[
\left| |E(S,T)| - \frac{r|S||T|}{n} \right| \leq \lambda r \sqrt{|S||T|}
\]
\end{lemma}

\begin{remark}[Spectral vs Vertex Expansion]
Spectral expansion implies vertex expansion. Specifically, Tanner's bound states that for any set $S \subset V$, $|N(S)| \ge \frac{|S|}{(1-\alpha)\lambda^2 + \alpha}$, where $\alpha = |S|/n$. If $\lambda$ is sufficiently small, the graph exhibits linear vertex expansion, which is essential for the distance properties in Theorem~\ref{thm:expander}.
\end{remark}

\subsection{Graph Separators}

\begin{definition}[Graph Separator]
A subset of vertices $S \subset V$ is a separator of a graph $G=(V,E)$ if the removal of $S$ partitions $V$ into two disjoint sets $A$ and $B$ such that $|A|, |B| \le \frac{2}{3}|V|$ and no edges connect $A$ and $B$.
\end{definition}

\begin{theorem}[Geometric Separator Theorem \cite{MillerTengThurstonVavasis}]\label{thm:geometric_separator}
Let $G$ be a graph embedded in $\mathbb{R}^2$ with density $\rho$ and interaction range (edge length) bounded by $L$. Then $G$ admits a separator of size $O(\sqrt{\rho} L \sqrt{|V|})$.
\end{theorem}

\section{Main Results}

\subsection{Tradeoff Bounds Under Combined Constraints}

\begin{theorem}[2D locality tradeoff (BPT, with range bookkeeping)]\label{thm:main}
Consider a 2D geometrically local stabilizer code on $n$ qubits with bounded density $\rho$ and interaction range $L$. Assume the stabilizer generators define geometrically local \emph{commuting constraints} in the sense of \cite{BPT}. Then there exists a constant
\[
C_{\mathrm{BPT}} = C_{\mathrm{BPT}}(\rho,L) > 0
\]
such that
\[
k\, d^2 \le C_{\mathrm{BPT}}\, n.
\]
Equivalently, writing $R=k/n$,
\[
d\,\sqrt{R} \le \sqrt{C_{\mathrm{BPT}}}.
\]
In particular, for fixed $\rho=O(1)$ and constant interaction range $L=O(1)$, one has $k d^2 = O(n)$.
\end{theorem}

\begin{proof}
The BPT theorem \cite{BPT} proves $k d^2 \le c\,n$ for 2D geometrically local commuting constraints with constant interaction range, where $c$ depends only on the locality parameters (range, maximum term size) and the local Hilbert space dimension. To incorporate a general interaction range $L$, one coarse-grains the plane into square blocks of side length $\Theta(L)$. Under bounded density $\rho$, each block contains at most $O(\rho L^2)$ qubits and can be treated as a single effective site. After coarse-graining, each constraint has constant range in the block lattice (range $O(1)$), at the expense of increasing the local Hilbert space dimension to $2^{O(\rho L^2)}$.

Applying BPT to the coarse-grained model yields $k d^2 \le C_{\mathrm{BPT}}(\rho,L)\, n$ for some constant $C_{\mathrm{BPT}}(\rho,L)$ determined by the resulting constant-range locality and effective on-site dimension. We do not attempt to optimize this dependence; the point is that for fixed $\rho$ and constant $L$, $C_{\mathrm{BPT}}$ is $O(1)$, recovering the usual 2D asymptotic tradeoff.
\end{proof}

\begin{corollary}
For 2D geometrically local codes with constant interaction range $L=O(1)$ and bounded density $\rho=O(1)$,
\[
d \le O\left(\sqrt{\frac{n}{k}}\right) \qquad\text{and}\qquad k d^2 \le O(n),
\]
with constants depending only on the locality model as in Theorem~\ref{thm:main}.
\end{corollary}


\subsection{Expansion and Linear Distance (Classical)}

\begin{theorem}[Expander-Code Distance from Unique-Neighbor Expansion]\label{thm:expander}
Let $C\subseteq \{0,1\}^n$ be a \emph{classical} linear code with parity-check matrix $H$ whose Tanner graph is $(d_v,d_c)$-regular (variable degree $d_v$, check degree $d_c$). Suppose there exists a constant $\alpha>0$ such that every set $S$ of variable nodes with $|S|\le \alpha n$ satisfies
\[
|N(S)| > \frac{d_v}{2}|S|.
\]
Then the (classical) minimum distance of $C$ is greater than $\alpha n$.
\end{theorem}

\begin{proof}
Let $x\in C$ be a nonzero codeword with support $S\subseteq [n]$. Since $Hx=0$, every check node in $N(S)$ is incident to an even number of variables from $S$. In particular, there can be no check node in $N(S)$ with exactly one neighbor in $S$.

Assume for contradiction that $|S|\le \alpha n$. If every check in $N(S)$ has at least two neighbors in $S$, then counting edges between $S$ and $N(S)$ gives
\[
d_v |S| = |E(S,N(S))| \ge 2|N(S)|,
\]
so $|N(S)|\le (d_v/2)|S|$, contradicting the hypothesis. Therefore no nonzero codeword can have weight $\le \alpha n$, and the classical distance satisfies $d>\alpha n$.
\end{proof}

\begin{remark}
For a CSS code specified by $(H_X,H_Z)$, expansion conditions like Theorem~\ref{thm:expander} can lower bound the distances of the \emph{underlying classical} codes defined by $H_X$ and $H_Z$. Translating this into a lower bound on the \emph{quantum} distance additionally requires controlling the quotient by the stabilizer group (e.g., excluding low-weight vectors in $\ker(H_X)$ that lie outside $\operatorname{rowspan}(H_Z)$).
\end{remark}

\subsection{The Locality-Expansion Tension}

The geometric separator theorem (Theorem~\ref{thm:geometric_separator}) imposes fundamental limitations on the expansion properties of geometrically local graphs, creating a tension with the requirements for good LDPC codes.

\begin{proposition}[2D locality implies vanishing expansion]\label{prop:local_not_expander}
Let $G=(V,E)$ be a graph with $n=|V|$ and maximum degree $\Delta$ that admits a geometric embedding in $\mathbb{R}^2$ with bounded density $\rho$ and interaction range (edge length) at most $L$. Then $G$ has a vertex separator of size $O(\sqrt{\rho}\,L\sqrt{n})$, and consequently the (edge) conductance satisfies
\[
\phi(G) = O\left(\frac{\sqrt{\rho}\,L\Delta}{\sqrt{n}}\right).
\]
In particular, for constant $\rho,L,\Delta$, no such family can have constant conductance (and hence cannot form a family of constant-degree expanders).
\end{proposition}

\begin{proof}
The separator statement is exactly Theorem~\ref{thm:geometric_separator}. Let $S$ be such a separator with $|S|=O(\sqrt{\rho}\,L\sqrt{n})$, and let $A,B$ be the two components after removing $S$ with $|A|,|B|\le 2n/3$. Without loss of generality, assume $|A|\ge n/3$.

All edges from $A$ to $V\setminus A$ must touch the separator $S$, so the number of boundary edges satisfies $|\partial_E A| \le \Delta |S| = O(\sqrt{\rho}\,L\Delta\sqrt{n})$. The conductance is
\[
\phi(G) \;=\; \min_{\emptyset\ne U\subseteq V,\,|U|\le n/2} \frac{|\partial_E U|}{\Delta |U|}
\;\le\; \frac{|\partial_E A|}{\Delta |A|}
\;=\; O\left(\frac{\sqrt{\rho}\,L\Delta}{\sqrt{n}}\right),
\]
as claimed.
\end{proof}

\section{Tightness and Examples}

\subsection{Toric Codes}
\begin{example}[Tight Construction for Local Codes]
Consider the standard Toric code defined on an $\ell \times \ell$ lattice on a torus.
The code parameters are:
\begin{itemize}
    \item Physical qubits $n = 2\ell^2$
    \item Logical qubits $k = 2$
    \item Distance $d = \ell = \sqrt{n/2}$
    \item Rate $R = k/n = 1/\ell^2$
\end{itemize}
This construction is strictly geometrically local in 2D with interaction range $L=O(1)$ and check weight $w=4$.
Applying the BPT tradeoff (Theorem~\ref{thm:main}), we obtain $k d^2 = 2\ell^2 = \Theta(n)$, i.e., the toric code is tight up to constants for 2D local commuting constraints.
\end{example}

\subsection{Breaking the Barrier: Non-Local Codes}
Recent breakthroughs by Panteleev and Kalachev \cite{PK} and others have produced "good" quantum LDPC codes with parameters $[[N, \Theta(N), \Theta(N)]]$. These codes achieve linear distance $d = \Theta(N)$, which is impossible in 2D under Theorem \ref{thm:main} when the interaction range satisfies $L=O(1)$, the maximum stabilizer weight satisfies $w=O(1)$, and the rate satisfies $R=\Theta(1)$ (since then Theorem \ref{thm:main} implies $d=O(1)$).
This is consistent with our results because these asymptotically good codes \textbf{cannot} be embedded in any fixed-dimensional lattice with local checks. Their interaction graphs have expansion properties that forbid low-dimensional embeddings (the geometric separator condition of Theorem~\ref{thm:geometric_separator} fails).
Our Theorem \ref{thm:main} essentially quantifies the "price of locality": to achieve linear distance, one must abandon fixed-dimensional geometric locality.
At a high level, strong expansion is a key ingredient enabling linear-distance behavior in LDPC constructions; Theorem~\ref{thm:expander} records a standard expander-code implication on the classical side.

More precisely, geometric locality (bounded density and bounded interaction range) implies sublinear separators in any bounded-degree locality graph derived from the geometry; in 2D one expects separators of size $O(\sqrt{n})$ (up to model-dependent constants and coarse-graining at scale $L$), whereas constant-degree expanders necessarily have linear-size separators. This separator-versus-expansion tension is one mechanism behind why expander-based LDPC constructions are incompatible with fixed-dimensional geometric locality; see, e.g., \cite{LiptonTarjan,MillerTengThurstonVavasis}.


\section{Discussion and Open Problems}

Our results advance the understanding of quantum LDPC code limitations under realistic hardware constraints, bridging the gap between abstract expander-based constructions and geometrically constrained implementations. However, several significant questions remain open.

\begin{enumerate}
    \item \textbf{Lattice Geometry Dependence:} Can the dependence on the interaction range $L$ and the constant $C$ in Theorem~\ref{thm:main} be improved for specific lattice geometries, such as hyperbolic tilings, which naturally support better expansion properties than Euclidean grids?
    \item \textbf{Stoquastic Hamiltonians:} What is the optimal tradeoff when stoquasticity constraints are imposed on the Hamiltonian? Stoquastic Hamiltonians (those with non-positive off-diagonal matrix elements) are easier to simulate but may have more restrictive ground state properties relevant to code distance.
    \item \textbf{Probabilistic Bounds:} How do these bounds change under probabilistic constructions? While worst-case bounds provide hard limits, random instantiations of local codes might exhibit better average-case behavior, potentially allowing for looser practical constraints.
\end{enumerate}

\section{Conclusion}

In this work, we rephrased the 2D BPT locality tradeoff in a form that makes the dependence on interaction range and maximum stabilizer weight explicit, and we illustrated tightness via the toric code. We also highlighted why asymptotically good quantum LDPC codes must give up fixed-dimensional geometric locality, pointing to expansion as a recurring structural feature in such constructions.

\begin{thebibliography}{99}

\bibitem[BPT10]{BPT}
S.~Bravyi, D.~Poulin, and B.~Terhal.
\newblock Tradeoffs for reliable quantum information storage in 2D systems.
\newblock {\em Physical Review Letters}, 104(5):050503, 2010.

\bibitem[BE21]{PRXQuantum}
N.~Breuckmann and J.~Eberhardt.
\newblock Quantum low-density parity-check codes.
\newblock {\em PRX Quantum}, 2:040101, 2021.

\bibitem[PK22]{PK}
P.~Panteleev and G.~Kalachev.
\newblock Asymptotically good quantum and locally testable classical LDPC codes.
\newblock In {\em Proceedings of the 54th Annual ACM SIGACT Symposium on Theory of Computing (STOC)}, pages 375--388, 2022.

\bibitem[LT79]{LiptonTarjan}
R.~J. Lipton and R.~E. Tarjan.
\newblock A separator theorem for planar graphs.
\newblock {\em SIAM Journal on Applied Mathematics}, 36(2):177--189, 1979.

\bibitem[MTTV97]{MillerTengThurstonVavasis}
G.~L. Miller, S.-H. Teng, W.~Thurston, and S.~A. Vavasis.
\newblock Separators for sphere-packings and nearest neighbor graphs.
\newblock {\em Journal of the ACM}, 44(1):1--29, 1997.

\end{thebibliography}

\end{document}
